\label{chap.10}
\section{Summary}
\label{Summary}
This thesis asked whether a 6G slicing controller allocates bandwidth better when it knows the demand of the next interval, and whether decentralised multi-agent reinforcement learning suits the task better than a single central agent. We built a pipeline of self-supervised learning, a temporal graph neural network and multi-agent reinforcement learning, and tested it in a unified slicing environment built on real Telecom Italia Milan traffic. All results come from a held-out week, with checkpoints and rule targets chosen on separate validation days.

\paragraph{Does knowing the next interval help?}
For the rule-based controller, yes, by a modest amount. Replacing the current demand by the TGNN forecast raises the reward by about 0.03 at every size from 9 to 100 cells, on all seven test days, and recovers about one fifth of the gap to perfect knowledge (Chapter~\ref{chap.8}). The forecaster itself beats persistence on all 10{,}000 cells, with mean squared error lower by 18\%, 33\% and 31\% for eMBB, URLLC and mMTC. For the learned controller the answer is not established: giving it the forecast changed its reward by less than the seed-to-seed variation.

\paragraph{Is multi-agent learning better than a single central agent?}
Yes, in two ways. A parameter-shared actor of 19{,}332 parameters works from 9 to 100 cells, while the central agent needs 114{,}468 to 604{,}048 parameters and its reward falls from about 0.43 to 0.27 (Chapter~\ref{chap.7}). With a long training budget, checkpoints chosen on validation days and seeds averaged, the learned controller also beats a simply tuned expert rule on every test day at 9, 25 and 100 cells, by 0.023 to 0.026 in reward. The off-policy controller does not clearly beat a richer rule that is tuned for each time of day and cell class: its difference to that rule is between +0.001 and +0.005 and its interval includes zero at 25 and 100 cells. The on-policy learners (IPPO and MAPPO, 12 runs) do beat that rule: their mean reward is 0.908 against 0.899, a gain of 0.009 $\pm$ 0.002 on 7 of 7 days at 25 cells, and every run is better on all seven days. They do not reach the rule that is given the true next demand (0.920).

\section{Contributions}
\label{Contributions}
\begin{enumerate}
    \item \textbf{A unified slicing environment.} Since no public dataset joins topology, per-slice traffic, quality of service, SLA outcomes and allocation actions, we built an environment from real Milan traffic and grid data, with latency calibrated on public 6G RAN telemetry and load-coupled interference. Each component is marked as real, calibrated or modelled.
    \item \textbf{A fair evaluation protocol for MARL.} Train, validation and test periods are separate, baselines are tuned on the training days, results are paired by day with intervals, and the strongest rule is part of the comparison.
    \item \textbf{A measurement of what prediction is worth.} The value of the forecast is stated as the reward gained and as the fraction of the gap to perfect knowledge, at four network sizes.
    \item \textbf{A city-scale forecasting finding.} A forecaster trained on the dense centre failed in the periphery, with error 10.8 times that of persistence, because it had learned the traffic level of the centre. Per-cell scaling, patch training and more training days removed the failure.
    \item \textbf{A clean negative result for SSL.} Pretraining did not beat strong raw-feature baselines on de-duplicated 5G-NIDD flows at any label budget.
    \item \textbf{A Digital Twin that checks decisions.} The approval gate recovers the reward lost to poor proposals, and a Simu5G sweep tests the surrogate against a packet-level simulator.
    \item \textbf{A dataset and a prototype.} The dataset has 393.6 million rows and 122 columns, six what-if scenarios and injected incidents, and the Slice Twin monitor provides a visual interface to the results.
\end{enumerate}
The work also documents corrections that changed early results: a label-ordered split and duplicate flows in 5G-NIDD, a slice-indexing error in the B5G environment, a city-wide forecast bias, and an earlier estimate of the value of the forecast that was reduced from 0.05 to 0.03 once the test period was extended to a week.

\section{Limitations}
\label{sec:limitations}
\paragraph{Statistical strength.}
The test is one week (seven days, which are not independent) in one city, with at most three seeds per learned controller and one seed at 100 cells. The intervals reported describe the variation across those days, not across weeks, cities or training runs. A stronger study would use ten or more seeds and more weeks.

\paragraph{Comparison with the best hand-built rule.}
The off-policy learned controller only matches a rule tuned per time of day and cell class; the PPO-family learners beat it by about 0.009 reward, a small margin, and were tested only at 25 cells with the default QoS curve. The case for learning also rests on removing the manual tuning and on its behaviour as the network grows. The central-agent baseline is also weak: its reward fell during training, so the result shows that this central design performs poorly, not that every centralised design would.

\paragraph{Fidelity of the surrogate model.}
QoS comes from a calibrated surrogate, and the telemetry used for calibration is itself simulated. The latency spread was fitted on the same SLA-against-load curve that is used for evaluation, so the drop in the mean gap from 0.086 to 0.021 measures fit and is not independent validation. The Simu5G sweep agrees with the surrogate on the order of delays but not on their shape: the packet-level delay is flat until about 93\% of capacity and then rises steeply, whereas the surrogate rises gradually from about half of capacity. The slice mapping, demand scale, interference coupling and reward weights are modelling assumptions.

\paragraph{Negative results.}
Self-supervised pretraining did not outperform raw features for 5G-NIDD classification or for the Milan forecast, so the contribution of the SSL stage to the pipeline is not demonstrated. The forecast has not been shown to help the learned controller. Over the whole city the static split beats the predictive controller because the rule starves quiet slices, and the city comparison covers a single day. The forecaster was trained once, without an ablation of its parts and without graph-free or published spatio-temporal baselines.

\paragraph{The Digital Twin claim.}
The twin trains controllers, evaluates them, generates scenarios and gates proposals. The gate uses the same QoS model as the environment in which it is tested, so how well it would predict outcomes on a real network is not shown. The controllers have not been validated against a live or packet-level network. The prototype is a monitor of one held-out day with scripted incidents.

\section{Future Work}
\label{sec:future-work}
\paragraph{Stronger statistics and baselines.}
Repeat the experiments with ten or more seeds and several test weeks, add the 100-cell runs with more seeds, and add a graph-free forecaster and published spatio-temporal models together with an ablation of the city forecaster.

\paragraph{Closing the simulator gap.}
Connect the project's own 19-gNodeB scenario to Simu5G, run controllers in closed loop with it and record latency, jitter and loss, which would also supply the latency and loss data that the B5G dataset lacks. Calibrate the surrogate on the measured delay curve, in particular the flat-then-cliff shape.

\paragraph{SSL with richer inputs.}
Test SSL on inputs where hand-crafted features are not available, such as raw radio measurements or longer sequences, and with an encoder that uses neighbouring cells and earlier intervals.

\paragraph{Forecast-aware learning.}
Investigate why the forecast does not help the learned controller, for example with a longer training budget, a design in which the forecast is part of the critic, or training on forecast errors.

\paragraph{Semantic slicing and O-RAN traces.}
Extend the reward with semantic criteria suggested by the literature, and check the surrogate against public O-RAN and testbed traces that contain measured PRB use, throughput and channel quality.

\paragraph{Deployment.}
Move the monitor and controllers from replay towards a closed loop with a simulator or testbed, and study decision latency inside the control interval, coordination between neighbouring cells and the cost of reconfiguration.

\section{Concluding Remarks}
Within these limits, the work shows three things. Predicting traffic before allocating bandwidth helps a rule-based controller. A parameter-shared multi-agent controller scales to larger networks where a central one does not, and it beats rules whose targets were tuned on the training days. And a forecaster has to be trained on the whole city if it is going to be used there. Two questions stay open: whether a learned controller can make real use of the forecast, and whether a packet-level simulator agrees with our surrogate beyond the ordering of delays.
